% ============================================================
% CAPITOLO 2 - LA TONALITA'
% ============================================================
\chapter{La tonalità}
\label{chap:tonalita}

Una tonalità è una scala con un centro, la tonica, intorno a cui si dispongono gli altri gradi. Questo capitolo dà i nomi dei gradi, distingue il modo maggiore dal minore e mostra come le tonalità sono collegate fra loro.

\section{Gradi della scala e loro nomi}
\label{sec:gradi_scala}

\begin{definizione}[Gradi della scala]
	\label{def:gradi_scala}
	I sette gradi della scala si indicano con numeri romani e hanno un nome:
	\begin{center}
	\begin{tabular}{cl}
	\toprule
	\textbf{Grado} & \textbf{Nome} \\
	\midrule
	I & tonica \\
	II & sopratonica \\
	III & mediante \\
	IV & sottodominante \\
	V & dominante \\
	VI & sopradominante \\
	VII & sensibile (nel minore naturale: sottotonica) \\
	\bottomrule
	\end{tabular}
	\end{center}
\end{definizione}

La tonica è il centro di gravità della scala. La dominante, quinta sopra la tonica, e la sottodominante, quinta sotto, sono i due gradi che le stanno intorno. La sensibile dista solo un semitono dalla tonica e tende a salire su di essa. Le funzioni che questi gradi svolgono nel discorso armonico si studiano nel Capitolo \ref{chap:funzioni_armoniche}.

\section{Modo maggiore e modo minore}
\label{sec:modi}

Il modo maggiore e il modo minore si distinguono per la successione di toni e semitoni. Conta soprattutto l'intervallo tra il I e il III grado: una terza maggiore dà il modo maggiore, una terza minore il modo minore.

Una seconda asimmetria riguarda il VII grado. Nel modo maggiore c'è una vera sensibile. Nel minore naturale il VII grado è un tono sotto la tonica e si chiama sottotonica. Per ottenere la tensione necessaria alla cadenza lo si alza di un semitono, e nasce così la \textbf{scala minore armonica}.

L'alzamento della sottotonica crea però, tra il VI grado e la sensibile, una \textbf{seconda aumentata}, che la melodia tende a evitare (Capitolo \ref{chap:errori_condotta}).

\section{Il circolo delle quinte}
\label{sec:circolo_quinte}

Le relazioni tra le tonalità si leggono sul \textbf{circolo delle quinte}. Ogni passo corrisponde a una quinta giusta ascendente, o a una quarta giusta discendente, e aggiunge un diesis (o toglie un bemolle) all'armatura di chiave.

\begin{figure}[H]
	\centering
	\begin{tikzpicture}[scale=3.2]
		% --- Cerchi concentrici ---
		\draw[thick, gray!80] (0,0) circle (1.5cm);
		\draw[thick, gray!80] (0,0) circle (1.1cm);

		% --- Tonalità maggiori (cerchio esterno) ---
		\node at (90:1.3cm) {\textbf{Do}};
		\node at (60:1.3cm) {\textbf{Sol}};
		\node at (30:1.3cm) {\textbf{Re}};
		\node at (0:1.3cm) {\textbf{La}};
		\node at (330:1.3cm) {\textbf{Mi}};
		\node at (300:1.3cm) {\textbf{Si}};
		\node at (270:1.3cm) {\textbf{\Fadiesis{}/\Solbemolle{}}};
		\node at (240:1.3cm) {\textbf{\Rebemolle{}}};
		\node at (210:1.3cm) {\textbf{\Labemolle{}}};
		\node at (180:1.3cm) {\textbf{\Mibemolle{}}};
		\node at (150:1.3cm) {\textbf{\Sibemolle{}}};
		\node at (120:1.3cm) {\textbf{Fa}};

		% --- Tonalità minori (cerchio interno) ---
		\node[scale=0.85] at (90:0.8cm) {la m};
		\node[scale=0.85] at (60:0.8cm) {mi m};
		\node[scale=0.85] at (30:0.8cm) {si m};
		\node[scale=0.85] at (0:0.8cm) {fa\sharp\ m};
		\node[scale=0.85] at (330:0.8cm) {do\sharp\ m};
		\node[scale=0.85] at (300:0.8cm) {sol\sharp\ m};
		\node[scale=0.85] at (270:0.8cm) {re\sharp/mi\flat\ m};
		\node[scale=0.85] at (240:0.8cm) {si\flat\ m};
		\node[scale=0.85] at (210:0.8cm) {fa m};
		\node[scale=0.85] at (180:0.8cm) {do m};
		\node[scale=0.85] at (150:0.8cm) {sol m};
		\node[scale=0.85] at (120:0.8cm) {re m};

		% --- Elenco esplicito delle alterazioni (anello più esterno) ---
		\node[gray, scale=0.7] at (90:1.75cm) {Nessuna};
		\node[gray, scale=0.7] at (60:1.75cm) {\Fadiesis{}};
		\node[gray, scale=0.7, align=center] at (30:1.75cm) {\Fadiesis{}, \Dodiesis{}};
		\node[gray, scale=0.7, align=center] at (0:1.75cm) {\Fadiesis{}, \Dodiesis{}, \Soldiesis{}};
		\node[gray, scale=0.7, align=center] at (330:1.75cm) {\Fadiesis{}, \Dodiesis{}, \Soldiesis{}, \Rediesis{}};
		\node[gray, scale=0.7, align=center] at (300:1.75cm) {\Fadiesis{}, \Dodiesis{}, \Soldiesis{}, \Rediesis{}, \Ladiesis{}};
		\node[gray, scale=0.65, align=center] at (270:1.85cm) {%
			$\substack{\text{\Fadiesis{}, \Dodiesis{}, \Soldiesis{}, \Rediesis{}, \Ladiesis{}, \Midiesis{}} \\ \text{\Sibemolle{}, \Mibemolle{}, \Labemolle{}, \Rebemolle{}, \Solbemolle{}, \Dobemolle{}}}$};
		\node[gray, scale=0.7, align=center] at (240:1.75cm) {\Sibemolle{}, \Mibemolle{}, \Labemolle{}, \Rebemolle{}, \Solbemolle{}};
		\node[gray, scale=0.7, align=center] at (210:1.75cm) {\Sibemolle{}, \Mibemolle{}, \Labemolle{}, \Rebemolle{}};
		\node[gray, scale=0.7, align=center] at (180:1.75cm) {\Sibemolle{}, \Mibemolle{}, \Labemolle{}};
		\node[gray, scale=0.7, align=center] at (150:1.75cm) {\Sibemolle{}, \Mibemolle{}};
		\node[gray, scale=0.7] at (120:1.75cm) {\Sibemolle{}};

		% --- Frecce indicatrici ---
		\draw[->, very thick, red!70, line width=1.2pt] (100:2.1cm) arc (100:20:2.1cm);
		\node[red!70, rotate=-30] at (60:2.35cm) {\textbf{QUINTE} $\rightarrow$};

		\draw[->, very thick, blue!70, line width=1.2pt] (140:2.1cm) arc (140:220:2.1cm);
		\node[blue!70, rotate=90] at (180:2.3cm) {$\leftarrow$ \textbf{QUARTE}};
	\end{tikzpicture}
	\caption{Il circolo delle quinte e delle quarte con elenco esplicito delle alterazioni.}
	\label{fig:circolo_quinte_migliorato_compatto}
\end{figure}

Le tonalità vicine sul circolo condividono più note e più armonici, e l'orecchio le sente più vicine. Questa vicinanza sarà il criterio per scegliere le tonalità di arrivo di una modulazione (Capitolo \ref{chap:modulazione}). Una trattazione matematica del circolo, con l'aritmetica modulare, è nell'Approfondimento \ref{math:aritmetica_modulare}.
